\chapter{Analysis/Design of the pipelined architecture of the system}

The architecture of the proposed maritime surveillance system is designed in the form of a multi-stage, modular pipeline, where each stage serves a specific purpose, receiving enhanced data from the previous one and sending it to the next one, thereby providing a process clarity, as well as a real-time scalability. The system feeds two sets of data into the system, a continuous video stream from the drone, CCTV or satellite camera and an AIS data stream which includes vessel identifiers, positional coordinates, speed and course data. These are processed separately, then integrated in a sensor fusion layer that combines the spatial information of video-based tracking with the navigational information of the AIS telemetry. The video pathway sequentially detects ships using YOLOv8, tracks identities without loss using DeepOCSORT, and processes the data through homographic pixel-to-coordinate transformation to build the geo-referencing path, whereas the AIS pathway applies frame-rate interpolation to then match with the video path using Haversine and Hungarian assignment algorithms. Enriched tracks, which include MMSI, speed, course, and ship type, are then fused with noisy video detections using a Kalman filter that reconciles the higher-frequency, but less precise, video stream with the more-sparse but more precise AIS observations to obtain a combined state estimate for the geographic position and velocity of each vessel. The resulting feature sequences (latitude, longitude, speed over ground, and trigonometric course encodings) are passed through an LSTM path predictor which predicts the future position of the vessel over the next K timesteps, followed by a post-processing stage of applying a land mask filter to ensure predicted paths remain in physically valid water areas.

\section{Details about the system architecture}
\subsection{Specifications for the design}
The maritime surveillance system being proposed is expected to meet the following technical requirements. The system can receive continuous video input with a minimum resolution of 1280×720, and processes the video at a goal frame rate of 30 frames per second, which ensures real-time video surveillance throughput. The YOLOv8 detection module is set to be a trade-off between detection sensitivity and suppressing false positives, with a confidence threshold of 0.5 and an Intersection over Union (IoU) threshold of 0.45 for Non-Maximum Suppression. The DeepOCSORT tracker is instantiated with the default values of the maximum age of 30 frames for the retention of the lost track, minimum hit count of 3 for the confirmation of a track, and cosine similarity as the distance measure used for appearance-based Re-Identification. The prediction module that is used by LSTM takes 30 timesteps as input and predicts the next K timesteps for the vessel's location. In the geospatial mapping module, the pixel coordinates are transformed into geographic latitude-longitude pairs with the aid of homographic transformation matrices generated with the camera calibration parameters. For real-time performance, NVIDIA CUDA acceleration is required and a minimum of 6GB VRAM is required. All processing and video ingestion operations are performed at the frame-level and the system is built using PyTorch as the main deep learning framework and OpenCV to process the frames.

Seeing as the main objective of any design is to gauge how the software will function, it is essential to perform some pre-analysis work beforehand.Considering the primary goal of any design is to analyze how the software will operate, it is necessary to undertake some pre-analysis work prior to the design.

A comprehensive pre-analysis was performed prior to the design of the system that involved reviewing the data sets, selecting a suitable model and algorithmic benchmarking so each selected component was well matched to the maritime environment.

\subsection{\textbf{Dataset Analysis}}
The quality of the annotations, diversity of the environment, and the proportion of vessel classes in the maritime video data were assessed. A variety of sea states, illumination conditions, vessel scales, and occlusion scenarios were selected as a priority to be included in the datasets, so as to allow the trained detection model to adapt well to the various scenarios of practical surveillance. YOLO format frame extraction and annotation were done, and bounding boxes were annotated for each instance of a vessel for training, validation and test sets.

\subsection{\textbf{Detection Model Selection}}
After comparing various YOLO models, YOLOv8 is chosen as the detection backbone.Compared to YOLOv5, YOLOv7 and YOLOv7tiny, YOLOv8 is selected as the detection backbone. The anchor-free architecture, optimized feature pyramid network, and enhanced performance in detecting small objects made YOLOv8 the best option for maritime ship detection, particularly in terms of small object detection capabilities, which is crucial for identifying ships in the distance.

\subsection{\textbf{Tracking Algorithm Selection}}
After experiments, DeepOCSORT is chosen to be the best algorithm to preserve identities in dense, occluded maritime scenes among the two algorithms.Considering the performance in preserving identities, DeepOCSORT is compared and selected as the best algorithm for the dense maritime scenes with occlusion. When comparing the resistance to identity switching using DeepOCSORT with OSNet-based Re-Identification to other motion-based trackers, such as motion-only trackers, DeepOCSORT is shown to be the best tracker for persistent monitoring of vessels.

\subsection{\textbf{Prediction Model Selection}}
Instead, the trajectory prediction was done using LSTM networks, rather than Kalman Filters or GRU variants. Kalman Filters are a computational efficient linear motion estimation but are unable to represent the non-linear manoeuvres found in the movement of maritime vessels that have curved movements. Unlike LSTMs, which try to learn complex temporal dependencies from the historical position sequences, the current approach cannot learn such dependencies. A kinematic linear regression fallback is implemented for data with less than 30 steps in the trajectory history to be used as input for the LSTM.

\subsection{\textbf{Sensor Fusion Analysis}}
Identified a key design opportunity for the complementary nature of video tracking with AIS telemetry. Video is good at observing high frequency, but it is noisy and subject to degradation due to environmental conditions, while AIS is inaccurate in its observations, but is sparse. Kalman filter fusion strategy was used to take advantage of both sources and combine their respective individual strengths to get a single coherent and stable estimate of the state of the vessel.



\section{Summary}
A detailed description of the design and methodology used in the proposed maritime surveillance system were presented in this chapter. From detection confidence thresholds to tracker initialisation parameters, to LSTM input window lengths to GPU requirements, the specifics of each module were set to define the boundaries of operation of the system. The pre-analysis results were presented, which helped to choose the YOLOv8, DeepOCSORT, LSTM, and Kalman filter fusion algorithms as the main components.The main components of the algorithms, such as YOLOv8, DeepOCSORT, LSTM, and Kalman filter fusion, were selected after a detailed pre-analysis, where each design decision was substantiated with comparative analysis with other algorithms. Then the two-phase pipeline architecture was developed to detail the flow of data from ingestion of raw video and AIS stream to detection, identity-preserving tracking, geo-referencing, sensor fusion, feature sequence construction to trajectory prediction, output generation with the constraint of the land-mask. The mathematical formulations were presented to give a rigorous theoretical background of the behavior of the mathematical system involved in each stage of the system, such as bounding box regression, IoU computation, Cosine similarity matching of ReID, Haversine distance calculation, Kalman filter state estimation and LSTM gating equations. Lastly, the experimental methods used to prepare the datasets, train the models and test the pipeline with the integrated system were presented, outlining the process of how the system is tested. The rest of this chapter and the subsequent chapter implement the design, which is fully described in this chapter.